% !TeX root = ../教材题库_全部条目_转写工作稿.tex
% Source: XJ-B1-C03-EX-P066-Y1436
% 第3章 / 3.1 函数 / 教材 p.66、67 / PDF 第73、74页
% OCR未校对：公式与文字均可能有误；未完成数学TeX转写；禁止直接发给学生
% 原题及相邻原解截图（核对用）：../题目原图/XJ-B1-C03-EX-P066-Y1436-01.webp, ../题目原图/XJ-B1-C03-EX-P066-Y1436-02.webp
% 本文件只在转写完成后作为题干使用；OCR 条目尚不是可靠题干。
\begin{bankitem}{3.1 例题 例题2}{第3章 / 3.1 函数 / 教材 p.66、67 / PDF 第73、74页}{XJ-B1-C03-EX-P066-Y1436}{OCR未校对：公式与文字均可能有误；未完成数学TeX转写；禁止直接发给学生}
\begin{ocrdraft}
[教材p.66 / PDF第73页]
\par
0二0]下列函数中哪个与函数3，一云eR）是同一个函数？
\par
（2）3，一於
\par
（3）
\par
Y\ocrglyph{00AF}
\par
解
\par
（1）函数3，一（x）2一()e[0，
\par
十。。）），它与函数3，一(x\ensuremath{\text{[OCR:EUR]}}R）的对应关系
\par
[教材p.67 / PDF第74页]
\par
相同，但定义域不相同，所以这个函数与函数3-工(xeR）不是同一个函数．
\par
（2）函数3，一=x(xeR)，它与函数3-x(xeR）不仅对应关系相同，而且
\par
定义域也相同，所以这个函数与函数3，一x(xeR）是同一个函数．
\par
／2
\par
（3）函数3，一
\par
='（'e（一亠，0）IJ（0，十亠）），它与函数y=x(x\ensuremath{\text{[OCR:EUR]}}R)的对应
\par
关系相同，但定义域不相同，所以这个函数与函数3-工(xeR）不是同一个函数．
\end{ocrdraft}
\end{bankitem}
